\section{Discussion}
\label{sec:discussion}

\subsection{Why PSO Did Not Beat Random Search}
\label{subsec:why}

In \genone the defaults appear to sit in a broad, flat region of good settings. They beat most proxy evaluations of both searches and were only 0.019 (PSO) and 0.023 (random search) behind the two winners (Section~\ref{subsec:res-search}). In the proxy check they ranked second of eight at both training lengths, and after 40 epochs the four best settings were within 0.026 of each other (Table~\ref{tab:proxy}).

Inside that region the proxy could tell poor settings from good ones ($\rho = 0.976$) but could not pick a winner. Its \genone winner, the random-search setting, had the lowest validation score after 100 epochs (Table~\ref{tab:gen1-test}). A short proxy may favor settings that learn fast, and taking the maximum of 32 single-run scores favors scores that were lucky. We did not measure either effect.

The budget was also small: 32 evaluations in six dimensions, and PSO spends eight of them on its random initial swarm. Bergstra and Bengio~\cite{bergstra2012random} report that for most data sets only a few hyperparameters really matter, and that random search finds models as good as or better than grid search at the same budget. If one or two of our dimensions dominate, 32 independent samples already cover them well and leave little for a guided search to exploit. All four search winners chose an initial learning rate below the default, which fits this view, although we did not measure the importance of each dimension.

The swarm also converged early. This raised its mean proxy score above that of random search but not its best score (Section~\ref{subsec:res-search}). In \gentwo the best particle was the last one of the random initial swarm (Fig.~\ref{fig:convergence}), so PSO in effect offered the best of eight random samples against the best of 32. With only three update steps, a global-best swarm pulls its later evaluations toward the best point of round~1. That improves the average, but only the maximum is used.

Why tuning helped only in \gentwo is less clear to us. Our guess is that the new data (real-world photographs, 640~px, a seventh class) moved the good region away from the defaults, so that any reasonable search gained about 0.02 \mapfifty on the mixed and studio test sets. We have no \gentwo proxy run of the defaults, so this is untested. Each arm was also trained only once, and part of the gain could be training noise. Repeating the proxy check and the full trainings for \gentwo would separate the two explanations (Section~\ref{subsec:fw-opt}).

\subsection{What This Means for Tuning Claims}
\label{subsec:implications}

Gap~3 (Section~\ref{subsec:gaps}) says that metaheuristic tuning in this area is compared with other metaheuristics or with nothing, using single test numbers. Our results show what that can hide. Had we compared PSO with the defaults alone, its significant \gentwo gain of $+0.020$ \mapfifty would have been credited to the swarm. Random search gained $+0.017$ at the same budget and the two tuned arms cannot be told apart (Table~\ref{tab:gen2-test}), so the gain comes from searching at all. Single numbers would have misled as well. The \genone point estimates rank the defaults first and PSO last, and \gentwo validation ranked random search ahead of PSO, yet neither ordering is significant on test. The opposite case also occurs: PSO is significantly better than random search on India by $+0.003$, on scores near 0.02, which shows that a significant difference still has to be large enough to matter.

From this we suggest a minimum protocol for a claim that a metaheuristic tunes a waste model better: (1) include the defaults and random search at the same number of evaluations, and report the time (Section~\ref{subsec:res-cost}); (2) check any cheap proxy against longer training before the search, with a threshold fixed in advance; (3) choose the final setting of each arm on validation data and score the test split once; (4) report paired confidence intervals on the differences between arms for each test domain, and judge them against a difference that would matter in practice; and (5) repeat trainings and searches with several seeds. Our study meets the first four and not the fifth. For an equal-budget comparison like that of Sayed et al.~\cite{sayed2024intelligent}, adding a random-search arm would cost one more competitor.

\subsection{Real-World and Indian Deployment}
\label{subsec:deployment}

For a smart bin or a municipal service the data matter far more than the optimizer. Every arm loses about half its studio score on real-world photographs and drops to 0.023 to 0.027 \mapfifty on Indian ones (Section~\ref{subsec:res-domain}). The largest tuning effect, about 0.02, is less than a tenth of the studio-to-real-world drop (about 0.28), and it did not appear on real-world photographs at all. The clearest evidence is Table~\ref{tab:g1g2}, which scores both generations on the same images. Real-world training data raised real-world \mapfifty from 0.038--0.096 (\genone) to 0.338--0.353 (\gentwo), with every \gentwo arm significantly above every \genone arm, whereas tuning made no significant real-world difference in either generation.

The price was a significant studio loss of 0.074 to 0.105. We see three likely causes, which our runs cannot separate: \gentwo kept only 41\% of the \genone studio training images; the same small YOLOv8n has to share its capacity with real-world scenes and a seventh class; and \gentwo trained at 640 instead of 416~px and chose its checkpoint on mixed validation data. Either way, local training data will help more than a better search, which is why Indian data come first in our future work (Section~\ref{subsec:fw-data}). Clutter is the second obstacle. The models label every object correctly on about 44\% of the single-object real-world photographs but on none of the 46 with three or more objects (Fig.~\ref{fig:best}). A bin that photographs one item at a time is a realistic target today, but sorting a mixed pile is not.

The DWSD result also carries a lesson about test sets. Its merged, mask-derived boxes penalize correct per-item detections (Section~\ref{subsec:res-domain}), so a test set has to be labeled in the same way as the training data before its score can be read as accuracy. Relabeling part of DWSD item by item would separate the labeling style from the model error.

Indian rules also constrain the classes. Rule~15(h) of the Solid Waste Management Rules, 2016~\cite{moefcc2016solid} requires green, white and black bins for bio-degradable, recyclable and other waste. BIODEGRADABLE maps to the green bin and the material classes, as far as their items are recyclable, to the white one. The black bin corresponds to OTHER, which is the weakest class of every arm (AP@50 0.050 to 0.063), so a bin built on these models would be least reliable exactly there.

\subsection{Limitations}
\label{subsec:limitations}

Each arm was trained once (seed 42) and each search was run once (PSO seed 42, random-search seed 123). The confidence intervals therefore cover test sampling only, and another seed could change the small differences between arms. We computed no $p$-values and no multiple-comparison corrections for the twelve pairwise comparisons per metric in \gentwo (three pairs on four test sets). Per-class AP has no interval, and the single BIODEGRADABLE box in the real-world test set makes its intervals wide and lopsided.

The proxy check used eight settings, was done for \genone only and used a 40-epoch reference instead of 100 epochs, so the validity of the proxy in \gentwo is assumed. The proxy scored last-epoch weights, whereas the final models used the best validation checkpoint (\texttt{best.pt}). The \gentwo Arm~C training was resumed after epoch 55 when a session froze; it kept the same weights and optimizer state but is not bit-identical to an uninterrupted run.

The class mapping of DWSD was inferred from the masks, with six of the fifteen mask values uncertain, and its boxes are connected regions and not single items. We searched only six training hyperparameters, and the budget was equal in evaluations, not in time. We did not run an external baseline such as the pipeline of Majchrowska et al.~\cite{majchrowska2022deep} on our data, so our scores cannot be ranked against prior work. Finally, the literature review was targeted and not systematic (Section~\ref{sec:related}), so the gaps of Section~\ref{subsec:gaps} hold only for the papers we reviewed.
